\chapter{JSON API, CORS és böngészős kliensintegráció}

\noindent\textbf{Tanulási út: Gyors út: API ág.} Most olvasd, ha JSON/API klienst vagy külön frontend origint használsz.\par\medskip

\rustbridgeblockhu{A typed skalárok, structok/modellek, collectionök, \code{Result} és újrahasznosítható business függvények ugyanazt az alkalmazásmagot adják, mint HTML route-oknál.}{A JSON input mód, zárt request séma, explicit public projection, content negotiation és CORS policy definiálja az API-határt.}{Ne építs második business réteget csak a JSON kedvéért. Ugyanazt a typed logikát használd, csak a reprezentáció és a boundary contract változik.}


\invariantblock{Biztonsági invariáns: a CORS nem authorization}{A CORS azt szabályozza, mely böngészős originekből indítható vagy olvasható cross-origin kérés. Nem hitelesíti a hívót, nem autorizál objektumot, és nem helyettesíti a normál request-életciklust.}

\diagnosticblockhu{Security diagnostic}{SEC-DATA-004}{Egy model explicit public projection nélkül kerülne kiadásra, ezért belső mezők véletlenül átléphetnék a response határt.}{Definiáld és add vissza a szándékolt public projectiont a belső model közvetlen szerializálása helyett.}

A Velran nem kényszerít választásra a klasszikus szerver-renderelt HTML és a JSON API között. Ugyanabban az alkalmazásban lehetnek \code{Html}-t visszaadó oldalak, formot feldolgozó actionök és typed JSON endpointok. Ez hasznos akkor is, ha csak egy kisebb JavaScript komponens kér adatot, és akkor is, ha külön frontend alkalmazás használja a Velran szervert API-ként.

A fejezet célja három működési modell tiszta szétválasztása:

\begin{enumerate}
  \item \textbf{same-origin HTML + API}: a böngésző ugyanarról az originről tölti az oldalt és hívja az API-t;
  \item \textbf{cross-origin publikus API}: külön origin hívja az API-t, de nincs session-cookie;
  \item \textbf{cross-origin session auth}: külön frontend origin használ cookie-alapú sessiont, ezért CORS és CSRF policy együtt szükséges.
\end{enumerate}

A legfontosabb alapelv: \textbf{a CORS nem autentikáció és nem CSRF-védelem}. A CORS böngészőoldali hozzáférési policy; az auth azt dönti el, ki a felhasználó, a CSRF pedig azt védi, hogy egy autentikált session ne legyen más originről jogosulatlanul felhasználható állapotmódosításra.

\section{Ugyanez az életciklus érvényes az API-kra is}

A JSON a reprezentációt változtatja meg, nem a trust modellt. Egy állapotmódosító API requestet ugyanabban a sorrendben érdemes végiggondolni:

\begin{lstlisting}
HTTP request
  -> parse: closed typed JSON schema
  -> validate: bounds/refinements
  -> authenticate: session or machine identity if required
  -> authorize: object/permission/tenant decision
  -> query / transaction
  -> response: explicit JSON projection or stable error envelope
\end{lstlisting}

A CORS e mellett az életciklus mellett helyezkedik el, nem valamelyik lépés helyett. Azt dönti el, hogy egy böngésző-origin a frontend kód számára hozzáférhetővé teheti-e a választ. Nem hoz létre identityt, objektum-authorityt, CSRF-biztonságot vagy adatbázis-integritást. Ha a CORS-t a core authorization láncon kívül tartod, elkerülhető a gyakori tervezési hiba: egy engedélyezett origint engedélyezett felhasználónak tekinteni.

\section{Az első typed JSON endpoint}

Egy JSON endpoint handlerének visszatérési típusa \code{Json}. A \code{json(...)} a támogatott runtime értéket JSON reprezentációvá alakítja.

\begin{lstlisting}[language=Velran]
#[page]
fn statusApi(ctx: PageContext) -> Result<Json, PageError> {
    let healthy = true;
    return Ok(json(healthy));
}

route statusApi GET "/api/status" public => statusApi;
\end{lstlisting}

A sikeres válasz reprezentációja JSON, például:

\Needspace{9\baselineskip}
\begin{lstlisting}
HTTP/1.1 200 OK
Content-Type: application/json; charset=utf-8

true
\end{lstlisting}

A \code{json(value)} nem tetszőleges objektum-serializáló. Scalar érték közvetlenül kódolható, de database/domain model publikus JSON boundaryn explicit \code{expose(...)} projectionön keresztül léphet át. Így egy új model field---különösen egy sensitive mező---nem jelenhet meg észrevétlenül egy meglévő API-ban. A response shape tehát tudatos szerződés, nem egy runtime objektum automatikus lenyomata.

\section{Adatbázisból JSON lista}

A typed SQL és a JSON response közvetlenül összekapcsolható. A query továbbra is pontos model shape-et ad vissza, a JSON réteg pedig ezt reprezentálja.

\begin{lstlisting}[language=Velran]
model Product {
    id: i64
    name: String
    price: i64
}

#[query]
fn listProducts(db: Db, offset: i64)
    -> Result<Vec<Product>, DbError> sql {
    SELECT id, name, price
    FROM products
    ORDER BY id
    LIMIT 100 OFFSET :offset
}

#[page]
fn productsApi(ctx: PageContext, db: Db,
                    offset: i64)
    -> Result<Json, PageError> {
    let products = listProducts(db, offset)?;
    return Ok(json(expose(products, id, name, price)));
}

route productsApi GET "/api/products"
    query offset<i64>
    validate offset range 0 1000000
    public => productsApi;
\end{lstlisting}

A fontos rész nem pusztán a JSON encoding. A reusable list query saját statikus row boundot hordoz; a kliens csak a bounded offsetet adja, nem az SQL cardinality authorityját. A route input typed és validált marad, a response pedig explicit public projectiont használ. Így egy később hozzáadott model field nem válik észrevétlenül az API részévé.

\Needspace{17\baselineskip}
\section{Typed JSON POST body}

JSON request bodyhoz a route \code{json} input módot deklarál.

\begin{lstlisting}[language=Velran]
#[action]
fn echoApi(ctx: ActionContext,
                  name: String,
                  age: i64,
                  active: bool)
    -> Result<Json, PageError> {
    return Ok(json(name));
}

route echoApi POST "/api/echo"
    json name<String> age<i64> active<bool>
    validate name length 1 100 age range 0 150
    auth user
    => echoApi;
\end{lstlisting}

Egy megfelelő kérés:

\begin{lstlisting}
POST /api/echo HTTP/1.1
Content-Type: application/json
Accept: application/json
X-CSRF-Token: <session CSRF token>

{"name":"Alice","age":42,"active":true}
\end{lstlisting}

A JSON input schema \textbf{zárt}. A runtime hibának tekinti többek között az ismeretlen vagy hiányzó kulcsot, a duplikált kulcsot, a rossz scalar típust, valamint a nem támogatott nested object, array, float vagy \code{null} inputot. Ezek request-contract hibák, ezért \code{400 Bad Request} választ adnak. Hibás vagy hiányzó \code{Content-Type} JSON route-on \code{415 Unsupported Media Type}. A mező- és request-méretlimitek ugyanúgy érvényesek, mint más request body esetén.

A \code{json}, \code{form} és \code{upload} body módok ugyanazon POST route-on kölcsönösen kizárják egymást, és ezt már a compiler ellenőrzi. JSON route \code{application/json}, egyszerű form route \code{application/x-www-form-urlencoded}, upload route pedig \code{multipart/form-data} requestet vár. Egy endpointnak legyen egyértelmű request contractja.

\section{Content negotiation és a 406 válasz}

A kliens az \code{Accept} headerrel jelezheti, milyen reprezentációt tud fogadni. JSON route-nál támogatott például:

\begin{lstlisting}
Accept: application/json
Accept: application/*
Accept: */*
\end{lstlisting}

Ha a handler JSON választ adna, de a kliens kizárólag \code{text/html} típust fogad el, a szerver \code{406 Not Acceptable} választ ad. Ez jobb annál, mintha a szerver csendben olyan reprezentációt küldene, amelyet a kliens explicit elutasított.

\section{API hibák}

JSON route-on a runtime által klasszifikált alkalmazáshibák ugyanazt az egyszerű envelope-ot használják: egyetlen stabil \code{error} kódot adnak vissza, miközben a részletes diagnosztika a szerveroldali logban marad. A kliens ne angol hibaüzenetek szövegét parse-olja, hanem a HTTP státuszt és ezt a kódot kezelje.

JSON-return route felismerése után a hibák JSON error envelope-ban érkeznek. Például:

\begin{lstlisting}
HTTP/1.1 401 Unauthorized
Content-Type: application/json; charset=utf-8

{"error":"unauthorized"}
\end{lstlisting}

Tipikus error kódok:

\begin{itemize}
  \item \code{bad_request};
  \item \code{unauthorized};
  \item \code{forbidden};
  \item \code{csrf_failed};
  \item \code{unsupported_media_type};
  \item \code{database_unavailable};
  \item \code{resource_limit};
  \item \code{internal_error}.
\end{itemize}

Ez különösen fontos böngészős kliensnél. Egy autentikált JSON route névtelen kérésre \code{401}-et ad, nem HTML login redirectet. A frontend így státuszkód alapján tud dönteni, nem kell HTML választ API hibaként értelmeznie.

\section{Same-origin: az alapértelmezett és legegyszerűbb modell}

Ha a HTML és az API ugyanazon scheme + host + port kombinációról érhető el, a böngésző szempontjából same-origin alkalmazásról beszélünk.

\begin{lstlisting}
https://app.example/
https://app.example/api/products
\end{lstlisting}

Ilyenkor általában nincs szükség CORS konfigurációra. Egy egyszerű kliens:

\begin{lstlisting}
const response = await fetch("/api/products?limit=20&offset=0", {
  headers: { Accept: "application/json" }
});

if (!response.ok) {
  throw new Error(`HTTP ${response.status}`);
}

const products = await response.json();
\end{lstlisting}

Ha nincs valódi igény külön frontend originre, maradj same-origin modellen: egy trust boundaryval és a CORS-bonyolultság nagy részével kevesebb.

\section{CORS: exact origin allowlist}

Cross-origin böngészős hozzáférés alapból nincs engedélyezve. A trusted server configban konkrét origineket kell felsorolni; wildcard origin nincs.

\begin{lstlisting}
[web]
cors_origins = ["https://frontend.example"]
cors_allow_credentials = false
\end{lstlisting}

A CLI megfelelője:

\begin{lstlisting}
velran-server --cors-origin https://frontend.example ...
\end{lstlisting}

Minden engedélyezett origint explicit sorolj fel; tetszőleges request \code{Origin} értéket soha ne tükrözz vissza.

\section{Preflight}

Bizonyos cross-origin requestek előtt a böngésző \code{OPTIONS} preflight kérést küld. A Velran csak akkor enged preflightot, ha:

\begin{itemize}
  \item az \code{Origin} konfigurált;
  \item a cél GET/POST route létezik;
  \item a kért request header támogatott.
\end{itemize}

Támogatott például a \code{Content-Type}, az \code{Accept} és az \code{X-CSRF-Token}; a preflight a route/origin policy része, nem általános allow-all OPTIONS endpoint.

\section{Credentialed CORS session-cookie-val}

Ha a frontend külön originről használ Velran session authot, credentialed CORS szükséges:

\begin{lstlisting}
[web]
cors_origins = ["https://frontend.example"]
cors_allow_credentials = true
\end{lstlisting}

Productionben ez csak HTTPS mellett engedélyezett. A session cookie ilyen konfigurációban \code{SameSite=None; Secure} tulajdonságot kap, mert a böngészőnek cross-site requestben is el kell küldenie.

A JavaScript oldalon a fetch kérésnek is explicit credential mód kell:

\begin{lstlisting}
const response = await fetch("https://api.example/api/account", {
  credentials: "include",
  headers: { Accept: "application/json" }
});
\end{lstlisting}

A fetch \code{credentials} beállítása csak a megfelelő cookie-k elküldését engedi; alkalmazás-jogosultságot nem ad.

\section{CSRF JSON API-n}

Cookie-alapú sessionnél a böngésző a session cookie-t automatikusan kezeli, ezért állapotmódosító JSON requestnél ugyanúgy CSRF-védelem kell, mint HTML formnál.

Form POST esetén a token body mező:

\begin{lstlisting}
_csrf=<token>
\end{lstlisting}

JSON POST esetén header:

\begin{lstlisting}
X-CSRF-Token: <token>
\end{lstlisting}

Külön frontend originhez készíthető autentikált token endpoint:

\begin{lstlisting}[language=Velran]
#[page]
fn csrfApi(ctx: PageContext) -> Result<Json, PageError> {
    return Ok(json(csrfToken));
}

route csrfApi GET "/api/csrf" auth user => csrfApi;
\end{lstlisting}

A token ugyanahhoz a sessionhöz tartozik; a CORS nem helyettesíti a CSRF-validációt.

\section{Teljes böngészős POST folyamat}

Külön session-auth frontendnél a böngésző session cookie-t és sessionhöz kötött CSRF tokent küld; a runtime az üzleti handler előtt ellenőrzi az origint, authenticationt, CSRF-et és typed inputot.

Példa kliens:

\begin{lstlisting}
const csrf = await fetch("https://api.example/api/csrf", {
  credentials: "include",
  headers: { Accept: "application/json" }
}).then(async r => {
  if (!r.ok) throw new Error(`CSRF HTTP ${r.status}`);
  return r.json();
});

const result = await fetch("https://api.example/api/echo", {
  method: "POST",
  credentials: "include",
  headers: {
    "Content-Type": "application/json",
    Accept: "application/json",
    "X-CSRF-Token": csrf
  },
  body: JSON.stringify({ name: "Alice", age: 42, active: true })
}).then(async r => {
  if (!r.ok) throw new Error(`API HTTP ${r.status}`);
  return r.json();
});
\end{lstlisting}

\section{Reverse proxy és origin}

Apache vagy Nginx mögött különösen fontos, hogy a public scheme/host és a trusted proxy konfiguráció összhangban legyen. Az alkalmazás ne a kliens által szabadon küldött forwarded headerekből találja ki, hogy milyen origin tekinthető biztonságosnak. A trusted proxy boundary alapját a telepítési fejezet mutatja be; a teljes \code{server.toml} kulcsreferencia a könyv végén található.

A CORS allowlistbe mindig a \textbf{böngésző által látott frontend origint} írd, például:

\begin{lstlisting}
https://frontend.example
\end{lstlisting}

nem pedig a belső proxy vagy container címet.

\section{Mit ne tegyél?}

\begin{itemize}
  \item Ne legyen permissive vagy reflektált CORS; explicit origineket használj.
  \item Ne kezeld a CORS-t authentication vagy CSRF-védelem helyett, és productionben ne küldj session credentialt HTTP-n.
  \item Egy POST route egy typed body módot használjon, az API input pedig legyen bounded.
  \item JSON kliens API-státuszokat, például 401/403-at kezeljen, ne HTML login redirectet.
\end{itemize}

\section{Melyik architektúrát válaszd?}

\begin{center}
\begin{tabularx}{\textwidth}{@{}L{0.27\textwidth}Y@{}}
\toprule
\textbf{Helyzet} & \textbf{Ajánlott modell} \\
\midrule
szerver-renderelt oldal + kisebb JS & same-origin HTML + JSON API \\
azonos domainen futó SPA & same-origin API, CORS nélkül \\
külön publikus frontend, nincs session & exact-origin CORS credentials nélkül \\
külön frontend + session auth & exact-origin credentialed CORS + HTTPS + CSRF \\
server-to-server kliens & CORS nem releváns; az auth és hálózati policy a lényeg \\
\bottomrule
\end{tabularx}
\end{center}

\section{Kipróbálható példa}

A repository teljes mintája: \path{examples/json-api/app.vrn}.

Indulj same-originból; külön frontend origint és CORS-t csak akkor adj hozzá, ha a deployment ezt ténylegesen megköveteli.

\section{Ellenőrzött companion példák}
A repository companion példák:
\begin{itemize}
  \item \path{examples/json-api/app.vrn} -- JSON API route-flow;
  \item \path{examples/typed-query-form/main.vrn} -- típusos query/form structok;
  \item \path{examples/json-typed-response/app.vrn} -- típusos JSON response-határ.
\end{itemize}
Teljes request/response példához ezeket használd az izolált könyvrészlet helyett.

\section{Gyakorlat: válaszd külön az API- és a CORS-tervet}
\paragraph{Gyakorlási mód: támogatott.} Használd az előszóban meghatározott támogatott módszert.

Írd le a JSON request/response contractot anélkül, hogy böngészőt említenél. Ezután külön döntsd el, szüksége van-e egy másik originről futó böngészőnek engedélyre a híváshoz. Így nem keveredik össze a CORS az authenticationnel vagy az API-autorizációval.

\section{Tipikus hiba: alapból permissive CORS-t kapcsolunk be}
A CORS böngészős origin-policy, nem általános hálózati security mechanizmus. Ahol ésszerű, preferáld a same-origin telepítést; ha kell cross-origin hozzáférés, az exact originöket és credential-viselkedést tudatosan sorold fel.

\section{Folyamatos mintaprojekt: Secure Notes}
Indulj az \path{examples/secure-notes/checkpoints/09-json-api/main.vrn} snapshotból és annak ellenőrzött JSON companion példáiból. Adj read-only JSON reprezentációt a hitelesített felhasználó jegyzeteiről. Ugyanazokat az autorizációs szabályokat tartsd meg, mint a HTML page esetén. CORS-t csak akkor adj hozzá, ha szándékosan külön frontend origint telepítesz; egyébként maradjon same-origin az API.
